\chapter{Diskusjon}
Løsningen har fått tydeligere grenser mellom normal kontroll, valgfri
statusvisning og feil som skal låse kjøring. A-gaten reduserer risikoen for
utilsiktet start. Motorens egen startkontroll gir et ekstra kontrollpunkt.
Den begrensede feilhandleren bevarer årsak og stopper nye forsøk, mens
OLED-feilsperren lar et valgfritt grensesnitt falle bort uten gjentatt
bussarbeid. Disse egenskapene er testet i faktisk applikasjonskode.

Jevnere fremoverkjøring kommer fra at gode enkeltmålinger ikke gir stopp
eller en ny startkommando for hver runde. Skanningens tidlige retur og
fjerningen av en duplisert kontrollrunde før sving reduserer unødvendig
stasjonært arbeid. Sidemåling under servosentrering beholder filtreringen
og utnytter ventetiden bedre. Likevel kan konservativ filtrering, brede korridorkrav
og en utløpt 2000 ms-kandidat føre til hyppige stopp. Dette må observeres
på bilen før parametre endres. Mindre venting er ikke automatisk bedre
dersom servoen eller sensoren ikke har stabilisert seg.

Viktigste begrensning er manglende fysisk integrasjon av dagens versjon.
Native fakes kan injisere tidspunkter og feil presist, men representerer
ikke akustisk krysstale, spenningsfall, støy eller motorens treghet.
SPARK-beviset styrker lyspolicyens logiske egenskaper; det sier ingenting
om en løs ledning eller om motorene faktisk stopper. De eldre sensor-
og GPIO-prøvene må heller ikke brukes som validering av et nytt bilde.

Det finnes heller ikke grunnlag for å kalle dette en ferdig verifisert
sanntidsløsning. B-stopp er programvareavhengig, sensorstråler er sparsomme,
fart er ikke målt og systemets WCET er ikke dokumentert. Uavhengig fysisk
motorfrakobling må være mulig i forsøksoppsettet. Fart, akseptabel bane og
reaksjonsfrist må fastsettes ut fra målingene før autonom demonstrasjon.

For videre utvikling gir tidsstemplede sensordata, én tydelig motor-eier
og karakterisert sensor-/busstid et bedre grunnlag enn å dele alt i flere
tasks uten en konkret gevinst. Lukket hastighetsregulering eller scanning
i bevegelse kan vurderes når det finnes fartstilbakemelding, målt
servotid og en konservativ stoppmodell. En faglig god vurdering forutsetter
at gruppen forklarer både det som virker og det resultatene ennå ikke støtter.
